\documentclass[11pt,a4paper,oneside]{amsbook}
\usepackage[T1]{fontenc}
\usepackage[utf8]{inputenc}
\usepackage{lmodern}
\usepackage[american]{babel}
\usepackage{amsmath,amssymb,amsthm,mathtools}
\usepackage[all,cmtip]{xy}
\usepackage{microtype}
\usepackage{enumitem}
\usepackage{xspace}
\usepackage{geometry}
\geometry{margin=1.1in}
\usepackage[hidelinks]{hyperref}

% SGA notation used in the translated body fragments.
\DeclareMathOperator{\Hom}{Hom}
\DeclareMathOperator{\Spec}{Spec}
\DeclareMathOperator{\trace}{tr}
\DeclareMathOperator{\gr}{gr}
\DeclareMathOperator{\Coker}{Coker}
\DeclareMathOperator{\prof}{prof}
\DeclareMathOperator{\coprof}{coprof}
\newcommand{\id}{\mathrm{id}}
\newcommand{\kres}{\mathit{k}}
\newcommand{\isomto}{\overset{\sim}{\longrightarrow}}
\newcommand{\isomfrom}{\overset{\sim}{\longleftarrow}}
\newcommand{\ZZ}{\mathbf{Z}}
\newcommand{\NN}{\mathbf{N}}
\newcommand{\QQ}{\mathbf{Q}}
\newcommand{\RR}{\mathbf{R}}
\newcommand{\CC}{\mathbf{C}}
\newcommand{\PP}{\mathbf{P}}
\let\cal\mathcal
\let\goth\mathfrak
\let\it\mathit
\def\to{\mathchoice{\longrightarrow}{\rightarrow}{\rightarrow}{\rightarrow}}
\def\from{\mathchoice{\longleftarrow}{\leftarrow}{\leftarrow}{\leftarrow}}
\def\mto{\mathchoice{\longmapsto}{\mapsto}{\mapsto}{\mapsto}}
\def\To{\mathchoice{\Longrightarrow}{\Rightarrow}{\Rightarrow}{\Rightarrow}}
\def\lto#1{\mathchoice{\xrightarrow{\scriptstyle~#1~}}{\stackrel{#1}{\longrightarrow}}{}{}}
\def\tbigwedge{\mathop{\textstyle\bigwedge}\limits}
\def\cf{{cf.}\kern.3em}
\def\Cf{{Cf.}\kern.3em}
\def\ie{{i.e.}\kern.3em}
\def\resp{resp.\xspace}
\def\ptbl{.\kern.2em }
\newcommand{\No}{no.~}
\newcommand{\no}{no.~}
\def\quoi{}
\let\rd=d
\newcommand{\red}{\mathrm{red}}
\def\pr{\mathrm{pr}}

% References outside this exposé retain their source keys and
% printed numbers. They become ordinary references when a label is added.
\makeatletter
\newcommand{\SourceRef}[2]{\@ifundefined{r@#1}{#2}{\ref{#1}}}
\makeatother

\addto\captionsamerican{\renewcommand{\chaptername}{Expos\'e}}
\renewcommand{\thechapter}{\Roman{chapter}}
\renewcommand{\thesection}{\arabic{section}}
\theoremstyle{plain}
\newtheorem{theorem}{Theorem}[section]
\newtheorem{proposition}[theorem]{Proposition}
\newtheorem{lemma}[theorem]{Lemma}
\newtheorem{corollary}[theorem]{Corollary}
\newtheorem{corollaries}[theorem]{Corollaries}
\theoremstyle{definition}
\newtheorem{definition}[theorem]{Definition}
\theoremstyle{remark}
\newtheorem{remark}[theorem]{Remark}
\newtheorem{remarks}[theorem]{Remarks}
\newtheorem*{remarkstar}{Remark}
\newtheorem*{remarksstar}{Remarks}
\newtheorem*{scholium}{Scholium}
\newenvironment{namedstatement}[1]
  {\par\addvspace{.5\baselineskip}\noindent\textsc{#1.}\ \itshape\ignorespaces}
  {\par\addvspace{.5\baselineskip}}
\numberwithin{equation}{section}

\title[SGA~1, Expos\'e II (English draft)]{Smooth morphisms: generalities, differential properties\\
{\normalsize SGA~1, Expos\'e II\\English draft}}
\author{A.\ Grothendieck}
\date{}
\hypersetup{
  pdftitle={SGA 1, Expose II: Smooth morphisms (English draft)},
  pdfauthor={A. Grothendieck},
  pdfsubject={Unofficial English translation of SGA 1, Expose II}
}

\begin{document}
\frontmatter
\maketitle

\chapter*{About this draft}
This is an unofficial English translation of Expos\'e~II,
``Morphismes lisses: g\'en\'eralit\'es, propri\'et\'es diff\'erentielles'',
from SGA~1. The source volume is
A.\ Grothendieck and M.\ Raynaud,
\textit{Rev\^etements \'etales et groupe fondamental} (SGA~1),
S\'eminaire de g\'eom\'etrie alg\'ebrique du Bois Marie, 1960--61.
It follows the slightly corrected SMF recomposition,
\href{https://arxiv.org/abs/math/0206203v2}{arXiv:math/0206203v2}.

This draft contains the opening convention, all five sections,
and the closing errata, including proofs and footnotes.
Statement numbering and mathematical notation follow the source,
including its two statements numbered~1.1 (a definition and a
proposition) and its two statements numbered~4.18 (a corollary
and a block of remarks).
References to material outside this draft retain their original
numbers. Apparent mathematical misprints in the source have been
retained; they are recorded in the accompanying README, separately
from the translated text. Scholarly proofreading remains outstanding.

The translator's contribution is licensed under
the MIT License.
The French original remains copyright of its authors and publishers.

\tableofcontents
\mainmatter
\setcounter{chapter}{1}
\chapter{Smooth morphisms: generalities, differential properties}
\label{II}

% original p. 29
Cross-references to Expos\'e~\SourceRef{I}{I} are indicated by~I. One
recalls that rings are noetherian, and preschemes locally noetherian.

\section{Generalities}
\label{II.1}

Let $Y$ be a prescheme, let $t_1,\dots,t_n$ be
indeterminates, one sets
\begin{equation}
\label{eq:II.1.1}
Y[t_1,\dots,t_n]=Y\otimes_\ZZ \ZZ[t_1,\dots,t_n]\quoi.
\end{equation}
Thus $Y[t_1,\dots,t_n]$ is a $Y$-scheme, affine over~$Y$,
defined by the quasi-coherent sheaf of algebras
$\cal{O}_Y[t_1,\dots,t_n]$. The datum of a section of this
prescheme over~$Y$ is therefore equivalent to the
datum of~$n$ sections of~$\cal{O}_Y$ (corresponding to the images
of the~$t_i$ under the corresponding homomorphism). If $Y'$ is
over~$Y$, one has
\begin{equation}
\label{eq:II.1.2}
Y[t_1,\dots,t_n]\times_Y Y'=Y'[t_1,\dots,t_n]\quoi,
\end{equation}
(which implies that the datum of a $Y$-morphism from~$Y'$
to $Y[t_1,\dots,t_n]$ is equivalent to the datum of~$n$
sections of~$\cal{O}_{Y'}$), on the other hand one has
\begin{equation}
\label{eq:II.1.3}
\big(Y[t_1,\dots,t_n]\big)[t_{n+1},\dots,t_m]=Y[t_1,\dots,t_m]\quoi,
\end{equation}
by virtue of the analogous formula for polynomial rings
over~$\ZZ$. The formula~\eqref{eq:II.1.2} implies that
$Y[t_1,\dots,t_n]$ varies functorially with~$Y$.

$Y[t_1,\dots,t_n]$ is of finite type and flat over~$Y$.

\begin{definition}
\label{II.1.1}
Let $f\colon X \to Y$ be a morphism, making~$X$ into a
$Y$-prescheme. One says that $f$ is
\emph{smooth}\footnote{Old terminology: $f$ is \emph{simple}
at~$x$, or $x$ is a \emph{simple} point for~$f$. This
terminology led to confusion in various contexts
(simple algebras, simple groups) and had to be
abandoned.}
at $x \in X$, or that $X$ is \emph{smooth over~$Y$ at~$x$}, if there exists
an integer $n \geq 0$, an open neighborhood $U$ of~$x$, and an
\'etale $Y$-morphism from~$U$ to $Y[t_1,\dots,t_n]$. One says
that~$f$
(\resp $X$)
is \emph{smooth} if it is smooth at every point of~$X$. An
algebra~$B$ over a ring~$A$ is said to be smooth at a prime
ideal $\goth{p}$ of~$B$, if $\Spec(B)$ is smooth over~$\Spec(A)$
at
the point~$\goth{p}$;
% original p. 30
$B$ is said to be smooth over~$A$ if it is smooth over~$A$ at every prime
ideal~$\goth{p}$ of~$B$. Finally, a local homomorphism $A \to B$
of local rings is said to be smooth (or $B$ is said to be smooth over~$A$)\footnote{It is then better to say, as in EGA IV 18.6.1, that $B$
is ``\emph{essentially smooth}'' over~$A$.}
if $B$ is a localization of an algebra of finite type $B_1$ smooth
over~$A$.
\end{definition}

One notes that the notion of smoothness of~$X$ over~$Y$ is local on~$X$
and on~$Y$; if $X$ is smooth over~$Y$, it is locally of finite type
over~$Y$.

% The corrected source repeats 1.1. Preserve it with a distinct PDF anchor.
\setcounter{theorem}{0}
\begingroup
\renewcommand{\theHtheorem}{\theHsection.\arabic{theorem}.bis}
\begin{proposition}
\label{prop:II.1.1}
The set of points $x$ of~$X$ at which $f$ is smooth is open.
\end{proposition}
\endgroup

This is trivial from the definition.

\begin{corollary}
\label{II.1.2}
If $B$ is smooth over~$A$ at~$\goth{p}$, then it is smooth over~$A$
at~$\goth{q}$ for every prime ideal~$\goth{q}$ of~$B$ contained
in~$\goth{p}$.
\end{corollary}

\ref{prop:II.1.1} also implies that the last two
definitions~\ref{II.1.1} coincide in their common
domain of existence.

\begin{proposition}
\label{II.1.3}
\textup{(i)}~An \'etale morphism, in particular an open
immersion, an identity morphism, is smooth. \textup{(ii)}~An extension
of the base in a smooth morphism gives a smooth
morphism. \textup{(iii)}~The composite of two smooth morphisms is
smooth.
\end{proposition}

(i)~is trivial from the definition, more precisely one has:

\begin{corollary}
\label{II.1.4}
\'etale $=$ quasi-finite $+$ smooth.
\end{corollary}

(ii)~follows at once from the analogous fact for
\'etale morphisms (I~\SourceRef{I.4.6}{4.6}) and for the projections $Y[t_1,\dots,t_n]
\to Y$ (\cf \eqref{eq:II.1.2}). For (iii), this follows
formally from the fact that it is true separately for
``\'etale'' (I~\SourceRef{I.4.6}{4.6}) and for projections of the type
$Y[t_1,\dots,t_n]\to Y$
(\cf \eqref{eq:II.1.3}), and from the two facts cited for~(ii):
Suppose $Y$ is smooth over~$Z$ and $X$ is smooth over~$Y$, let us prove that $X$ is
smooth over~$Z$; one may suppose $Y$ \'etale over~$Z[t_1,\dots,t_n]$
and $X$ \'etale over~$Y[s_1,\dots,s_m]$, the first
hypothesis therefore implies that $Y[s_1,\dots,s_m]$ is \'etale over~$Z[t_1,\dots,t_n][s_1,\dots,s_m] = Z[t_1,\dots,s_m]$, hence $X$ is
\'etale over~$Z[t_1,\dots,s_m]$, qed.

\begin{remark}
\label{II.1.5}
The integer $n$ appearing in def\ptbl \ref{II.1.1} is well
determined, for one sees
% original p. 31
at once that it is the dimension of the local ring of~$x$ in its
fiber $f^{-1}\big(f(x)\big)$. One calls it the ``relative dimension''
of~$X$ over~$Y$. It behaves additively for the composition of
morphisms.
\end{remark}

\section{Some criteria of smoothness of a morphism}
\label{II.2}

\begin{theorem}
\label{II.2.1}
Let $f\colon X\to Y$ be a morphism locally of finite type, let $x\in
X$ and $y=f(x)$. For $f$ to be smooth at $x$, it is necessary and sufficient
that \textup{(a)} $f$ be flat at $x$, and \textup{(b)} $f^{-1}(y)$ be smooth over~$\kres(y)$
at $x$.
\end{theorem}

The composite of two flat morphisms being flat, and
$Y[t_{1},\dots,t_{n}]\to Y$ being a flat morphism, one sees that
smooth implies flat; taking~\ref{II.1.3}~(ii) into account this proves
necessity. Suppose (a) and (b) are satisfied, let $V$ be
an affine neighborhood of~$y$ with ring~$A$, $U$ an affine neighborhood of
$x$ over~$V$, with ring $B$. Taking $U$ small enough, one may
suppose by (b) that there exists a $\kres(y)$-morphism \emph{\'etale}
\[
g\colon U|f^{-1}(y)\to\Spec k[t_{1},\dots,t_{n}]\quad (k=\kres(y))
\]
defined by $n$ sections $g_{i}$ of the structure sheaf of
$U|f^{-1}(y)$. One easily sees that one may suppose that the
$g_{i}$ (which a priori are elements of
$B\otimes_{A}k=BS^{-1}$, where $S=A-\goth{p}$, $\goth{p}$ the prime
ideal of~$A$ corresponding to $y$) come from sections of
the structure sheaf of
$U$,
hence that $g$ is induced by a morphism,
still denoted~$g$
\[
g\colon U\to Y[t_{1},\dots,t_{n}]
\]
(up to multiplying the $g_{i}$ by one and the same
nonzero element of~$k$). Now
$U$
is flat over~$Y$ by (a), the same is true of
$Y[t_{1},\dots,t_{n}]$, on the other hand $g$ induces an
\'etale morphism between the fibers over~$y$, hence $g$ is \'etale at
$x$ by~(I~\SourceRef{I.5.8}{5.8}), qed.

\begin{corollary}
\label{II.2.2}
Let $S$ be a prescheme, $f\colon X\to Y$ an $S$-morphism of
finite type, $Y$ being of finite type and flat over~$S$, $x\in X$, $s$
the projection of~$x$ on~$S$. For $f$ to be smooth at $x$, it is necessary
and sufficient that $X$ be flat (or again: smooth) over~$S$ at $x$, and
that the morphism $f_{s}\colon X_{s}\to Y_{s}$ induced on the fibers of
$s$ be smooth at $x$.
\end{corollary}

Only sufficiency requires a proof, and follows from
the criterion~\ref{II.2.1}, together with the
flatness criterion~(I~\SourceRef{I.5.9}{5.9}).

To
% original p. 32
state the following result, ``recall'' that a morphism
$f\colon X\to Y$ locally of finite type is said to be
\emph{equidimensional}
at the point $x\in X$ if (setting $y=f(x)$) one can find an open
neighborhood $U$ of~$x$, every component of which dominates a component of~$Y$
such that, for every $y'\in Y$, the irreducible components of
$f^{-1}(y')\cap U$ all have one and the same dimension
independent of~$y'$. It is enough, moreover, in this condition to
take for $y'$ the generic points of the irreducible
components of~$Y$ passing through $y$, \emph{and} the point $y$. If
for example $X$ and $Y$ are integral and $f$ is dominant, the condition
means that the components of the $f^{-1}(y)$ passing through $x$ have ``the
right'' dimension, \ie the dimension of the generic fiber
(recall that they are always $\ge$ the dimension of the
generic fiber). If $f$ is equidimensional at $x$, the
dimension of its fiber at $x$ being $n$, and if $g\colon U\to
Y'=Y[t_{1},\dots,t_{n}]$ is a $Y$-morphism of a neighborhood~$U$ of
$x$, inducing a morphism on the fibers of~$y$ which is quasi-finite
at $x$ (or again, which amounts to the same, if $g$ is quasi-finite
at~$x$), then one shows that every irreducible component of~$U$
passing through $x$ dominates an irreducible component of~$Y'$. Moreover by virtue of the ``normalization lemma'', such a $g$
always exists (and conversely, if there exists a quasi-finite $Y$-morphism
$g$ of an open neighborhood $U$ of~$x$ into a $Y$-scheme
of the form $Y'=Y[t_{1},\dots,t_{n}]$, such that every component of~$U$ passing through $x$ dominates a component of~$Y'$, then $f$ is
equidimensional at~$x$). This being so:

\begin{proposition}
\label{II.2.3}
Let $f\colon X\to Y$ be a morphism locally of finite type, $x$ a
point of~$X$, $y=f(x)$, one assumes $\cal{O}_{y}$ normal. For $f$
to be smooth at $x$, it is necessary and sufficient that $f$ be
equidimensional at $x$, and that $f^{-1}(y)$ be smooth over~$\kres(y)$
at $x$.
\end{proposition}

One sees at once from the definition that a smooth morphism is
equidimensional (N.B. a flat morphism of finite type is not
necessarily equidimensional at $x$, even if its fiber
at $x$ is irreducible). Let us prove the converse. Since
$f^{-1}(y)$ is smooth over~$\kres(y)$ at $x$, one may suppose (replacing $X$ if need be by a suitable neighborhood of~$x$) that there exists
a $Y$-morphism
\[
g\colon X\to Y[t_{1},\dots,t_{n}]=Y'
\]
inducing an \'etale morphism on the fibers of~$y$, and a fortiori
quasi-finite at $x$. Hence
% original p. 33
$g$ is unramified, and ($f$ being equidimensional at
$x$) the irreducible components of~$X$ passing through $x$ each dominate
a component of~$Y'$, a fortiori the homomorphism
$\cal{O}_{y'}\to\cal{O}_{x}$ deduced from~$g$ (where $y'=g(x)$) is
\emph{injective}. This homomorphism is moreover unramified, and
$\cal{O}_{y'}$ is normal since it is a localization of the ring
$\cal{O}_{y}[t_{1},\dots,t_{n}]$, which is normal since
$\cal{O}_{y}$ is so. Hence the homomorphism $\cal{O}_{y'}\to\cal{O}_{x}$
is \'etale (I~\SourceRef{I.9.5}{9.5}~(ii)).

\begin{remarks}
\label{II.2.4}
The preceding statement still holds upon replacing
the hypothesis that $\cal{O}_{y}$ is normal by the weaker
hypothesis that $Y$ is \emph{geometrically unibranch} at $y$,
(\cf I~\SourceRef{I.11}{11}) --- since (I~\SourceRef{I.9.5}{9.5}) holds under this
hypothesis. Let us take this opportunity to note at the same time
that if the residue field of a local integral domain $A$ is
algebraically closed, then analytically integral
(\ie $\widehat{A}$ is an integral domain) implies geometrically
unibranch, the converse being moreover true in any
category of ``good rings'', more precisely in a
category of rings stable under the usual operations, and
in which the completion of a normal local ring is normal
(a condition fulfilled, by virtue of the ``analytic normality
theorem'' of Zariski, in the category of affine
algebras and their localizations)\footnote{\Cf EGA IV~7.8.}.

``Recall'' finally in the present context the following result,
due to Hironaka\footnote{\Cf EGA IV 5.12.10.} which sometimes allows
one to ensure that $f^{-1}(y)$ is a reduced scheme, \ie that
it is also what many algebraic geometers
would abusively regard as the fiber (without multiplicity)
of~$f$ above~$x$ (namely $f^{-1}(y)_{\red}$):
\end{remarks}

\begin{proposition}
\label{II.2.5}
Let $f\colon X\to Y$ be a dominant morphism of finite type of
reduced preschemes, $y$ a point of~$Y$ such that
$\cal{O}_{y}$ is regular. One assumes that all the components
of~$f^{-1}(y)$ are of multiplicity~$1$
(\cf the definition below), and that $f^{-1}(y)_{\red}$ is
normal. Then $f^{-1}(y)$ is reduced hence normal, $X$ is normal
at all the points of~$f^{-1}(y)$, finally $X$ is flat over~$Y$ at all
the points of~$f^{-1}(y)$.
\end{proposition}
One
% original p. 34
says that a component $Z$ of~$f^{-1}(y)$ is of
\emph{multiplicity}~$1$
if, $x$ denoting the generic point of~$Z$, one has (i)
$\dim\cal{O}_{x}=\dim\cal{O}_{y}$ (\ie $Z$ is not an ``excess
component'', that is, is not ``of too large
dimension''; (ii) the maximal ideal of~$\cal{O}_{x}$ is generated
by the maximal ideal of~$\cal{O}_{y}$ (which a priori, by virtue of
the choice of~$x$, generates an ideal of definition
of~$\cal{O}_{x}$).

Taking~\ref{II.2.3} or~\ref{II.2.1} into account one thus finds:

\begin{corollary}
\label{II.2.6}
Let $f\colon X\to Y$ be a dominant morphism of finite type of
reduced preschemes, $y$ a point of~$Y$ such that
$\cal{O}_{y}$ is regular. For $f$ to be smooth at the points of
$X$ above~$y$, it is necessary and sufficient that the components of
$f^{-1}(y)$ be of multiplicities $1$, and that $f^{-1}(y)_{\red}$
be smooth over~$\kres(y)$.
\end{corollary}

This situation was considered especially in the past
when $Y$ was the spectrum of a discrete valuation ring
$A$, and was commonly designated by phrases
such as: ``if the reduction of~$X$ with respect to the given
valuation is nice''... Moreover, $X$ then designated a
closed subscheme (if one may say so) of a $\PP_{K}^{n}$ ($K$
being the field of fractions of~$A$) and for want of an
adequate language, the more intrinsic role of an object ``defined
over~$A$'' (and not merely over~$K$) hardly appeared.

\section{Permanence properties}
\label{II.3}

\begin{proposition}
\label{II.3.1}
Let $f\colon X\to Y$ be a morphism, let $x\in X$ and
$y=f(x)$. Suppose $f$ is smooth at~$x$. For $\cal{O}_{x}$ to be
reduced (\resp regular, \resp normal) it is necessary and sufficient that
$\cal{O}_{y}$ be so.
\end{proposition}

This statement is indeed known when $X$ is of the form
$Y[t_{1},\dots,t_{n}]$, and it is proved in (I, \no \SourceRef{I.9}{9})
for an \'etale morphism; the general case is deduced from this
at once thanks to definition~\ref{II.1.1}.

We do not spell out here the other permanence
properties, already resulting from flatness alone, or from
the fact that $X$ is locally quasi-finite and flat over a
$Y$-prescheme of the form $Y[t_{1},\dots,t_{n}]$ (or, as
% original p. 35
we shall say, that $X$ is
Cohen--Macaulay
over~$Y$). Let us only note that from this last fact
it follows that
\begin{equation}
\label{eq:II.3.1}
{\dim\cal{O}_{x}=\dim\cal{O}_{y}+n-d,\quad
\prof\cal{O}_{x}=\prof\cal{O}_{y}+n-d},
\end{equation}
where $n$ is the dimension of the fiber of~$f$ at $x$, and $d$ the
transcendence degree of~$\kres(x)$ over~$\kres(y)$, whence (setting
$\coprof=\dim-\prof$)\footnote{For these formulas, \cf EGA~IV 6.1
and~6.3.}
\begin{equation}
\label{eq:II.3.2}
{\coprof\cal{O}_{x}=\coprof\cal{O}_{y}}.
\end{equation}
It follows for example that $\cal{O}_{x}$ is Cohen--Macaulay
(\resp without embedded components) if and only if the same is
true of~$\cal{O}_{y}$.

\section{Differential properties of smooth morphisms}
\label{II.4}

To simplify, we shall restrict ourselves for the most part to
differential calculus of order~$1$, confining ourselves to brief
indications for higher order (where the results are just as simple).

For the definition of the sheaf of $1$-differentials
$\it{\Omega}^1_{X/Y}$ of a $Y$-prescheme~$X$,
\cf (I~\No\SourceRef{I.1}{1}). Suppose that $X$ and $Y$ are
$S$-preschemes, the structure morphism $f\colon X\to Y$
being an $S$-morphism. Then $f$ defines a homomorphism of
Modules (compatible with~$f$)
\begin{equation}
\label{eq:II.4.1}
f^*\colon \it{\Omega}^1_{Y/S}\to\it{\Omega}^1_{X/S}
\end{equation}
in other words, $\it{\Omega}^1_{X/S}$ is \emph{contravariant} in
the $S$-prescheme~$X$. Moreover~\eqref{eq:II.4.1}
is equivalent to a homomorphism of Modules on~$X$
\begin{equation*}
\label{eq:II.4.1bis}
\tag{4.1\textrm{bis}}
f^*\big(\it{\Omega}^1_{Y/S}\big)\to\it{\Omega}^1_{X/S}
\end{equation*}
also denoted by~$f^*$ for want of a better notation, and
which fits into a canonical exact sequence of homomorphisms of
Modules
\begin{equation}
\label{eq:II.4.2}
f^*\big(\it{\Omega}^1_{Y/S}\big)\to
\it{\Omega}^1_{X/S}\to\it{\Omega}^1_{X/Y}\to 0
\end{equation}
All these homomorphisms are defined by the condition of being of
local nature (which reduces one to the affine case) and of commuting
with the operators~$\rd$. The exactness of~\eqref{eq:II.4.2} is
classical and trivial, and transcribes in the affine case as the
exact sequence (corresponding to a homomorphism $B\to C$ of
$A$-algebras):
\begin{equation*}
\label{eq:II.4.2bis}
\tag{4.2\textrm{bis}}
\Omega^1_{B/A}\otimes_{B}C\to\Omega^1_{C/A}\to\Omega^1_{C/B}\to 0
\end{equation*}

\begin{lemma}
\label{II.4.1}
Let
% original p. 36
$f\colon X\to Y$ be a morphism of $S$-preschemes. If $f$
is unramified (\resp \'etale) then
$f^*\big(\it{\Omega}^1_{Y/S}\big)\to\it{\Omega}^1_{X/S}$
is surjective (\resp an isomorphism). The converse is true
in the ``unramified'' case, if $f$ is assumed locally of
finite type.
\end{lemma}

The unramified case follows from the exact
sequence~\eqref{eq:II.4.2} and from (I~\SourceRef{I.3.1}{3.1}), but can also be
seen directly as the \'etale case. Consider the diagram
\begin{equation*}
\xymatrix@C=1.5cm{
X \ar[r]^-{\Delta_{X/Y}} & X\times_{Y}X \ar[d] \ar[r] &
X\times_{S}X \ar[d] \\ & Y \ar[r]^-{\Delta_{Y/S}} & Y\times_{S}Y
}
\end{equation*}
in which $X\times_{Y}X$ is identified with the fibered product of~$Y$ and
$X\times_{S}X$ over~$Y\times_{S}Y$. Since $f$ is unramified,
$X\to X\times_{Y}X$ is an open immersion, hence the
``conormal'' sheaf of the composite immersion $\Delta_{X/S}$ of the
latter with $X\times_{Y}X\to X\times_{S}X$ is isomorphic to
the inverse image on~$X$ of the conormal sheaf for the immersion
$X\times_{Y}X\to X\times_{S}X$. On the other hand, $X\to Y$ being
\'etale hence flat, $X\times_{S}X\to Y\times_{S}Y$ is flat, hence the
conormal sheaf for the immersion $X\times_{Y}X\to X\times_{S}X$ is
isomorphic to the inverse image of the conormal sheaf for the immersion
$Y\to Y\times_{S}Y$, \ie the inverse image of~$\it{\Omega}^1_{Y/S}$. The
conclusion follows from this.

\begin{lemma}
\label{II.4.2}
Let $X=Y[t_{1},\dots,t_{n}]$, $Y$ being an
$S$-prescheme. Then the sequence of canonical homomorphisms
\[
0\to f^*\big(\it{\Omega}^1_{Y/S}\big)\to
\it{\Omega}^1_{X/S}\to\it{\Omega}^1_{X/Y}\to 0
\]
is exact and $\it{\Omega}^1_{X/Y}$ is free with basis the
$\rd_{X/Y}t_{i}$.
\end{lemma}

The verification (purely affine) is immediate. (N.B.\ one
already knows the exactness of~\eqref{eq:II.4.2}).

Combining these two statements and definition~\ref{II.1.1}, one
finds

\begin{theorem}
\label{II.4.3}
Let $f\colon X\to Y$ be a smooth morphism of $S$-preschemes,
then:
\begin{enumerate}
\item[(i)] The sequence of canonical homomorphisms
\[
0\to f^*\big(\it{\Omega}^1_{Y/S}\big)\to
\it{\Omega}^1_{X/S}\to\it{\Omega}^1_{X/Y}\to 0
\]
is exact.
\item[(ii)] $\it{\Omega}^1_{X/Y}$ is locally free, its rank $n$
at~$x$ is equal to the relative dimension of~$f$ at~$x$.
\end{enumerate}
\end{theorem}

\begin{corollary}
\label{II.4.4}
The homomorphism
% original p. 37
$f^*\big(\it{\Omega}^1_{Y/S}\big)\to\it{\Omega}^1_{X/S}$ is
injective, its image in $\it{\Omega}^1_{X/S}$ is locally a
direct factor.
\end{corollary}

Let $u\colon F\to G$ be a homomorphism of Modules on the
prescheme~$X$; one says that it is \emph{universally
injective}
at $x\in X$ if the homomorphism $F_{x}\to G_{x}$ of
$\cal{O}_{x}$-modules is injective, and remains so upon tensoring
with any $\cal{O}_{x}$-algebra (or, which comes to the same thing
moreover, with any $\cal{O}_{x}$-module). It suffices for example
that there exist an open neighborhood $U$ of~$x$ such that $u$ induces an
isomorphism of~$F|U$ onto a direct factor of~$G|U$; this condition
is also necessary when $F$ and $G$ are free (and of finite
type) in a neighborhood of~$x$; more precisely in this case
the following conditions are equivalent:
\begin{enumerate}
\item[(i)] $u$ is injective at~$x$ and $\Coker u$ free at~$x$;
\item[(ii)] There exists an open neighborhood $U$ of~$x$ such that $u$
induces an isomorphism of~$F|U$ onto a direct factor of~$G|U$;
\item[(iii)] $u$ is universally injective at~$x$;
\item[(iv)] the homomorphism
$F_{x}\otimes \kres(x)\to G_{x}\otimes \kres(x)$
on the restricted ``fibers'' induced by~$u$ is injective;
\item[(v)] The transposed homomorphism $\check{G}\to\check{F}$ is
surjective at the point~$x$ (or again, which comes to the
same thing, in a neighborhood of~$x$).
\end{enumerate}
(Circular proof, (iv)$\To$(v) follows from
Nakayama, on the other hand (v)$\To$(i)
since a locally free quotient sheaf is necessarily a
direct factor). Geometrically, the situation envisaged
means that $u$ corresponds to an isomorphism of the vector bundle whose
sheaf of sections is~$F$, onto a subbundle of the
analogous vector bundle defined by~$G$. Of course, it does
not suffice for this that $F\to G$ be injective.

\begin{corollary}
\label{II.4.5}
Let $f\colon X\to Y$ be a morphism of $S$-preschemes,
locally of finite type, $x\in X$, $y=f(x)$, $s$ the projection of~$x$
and~$y$ on~$S$. One assumes $Y$ is smooth at~$y$ over~$S$. Equivalent
conditions:
\begin{enumerate}
\item[(i)] $f$ is smooth at~$x$.
\item[(ii)] $X$ is smooth over $S$
at~$x$, and
$f^*\big(\it{\Omega}^1_{Y/S}\big)\to\it{\Omega}^1_{X/S}$ is
universally injective at~$x$, \ie it is an injective homomorphism at~$x$ and its cokernel $\it{\Omega}^1_{X/Y}$ is free
at~$x$.
\end{enumerate}
\end{corollary}

Necessity follows from~\ref{II.1.3}~(iii) and
from~\ref{II.4.3}~(i)\,(ii); let us prove sufficiency. Since the
$\rd g$ ($g\in\cal{O}_{x}$) generate the module
$\it{\Omega}^1_{X/Y}$ at~$x$, one can find
% original p. 38
$g_{i}$ ($1\le i\le n$) such that the images of the $\rd g_{i}$ in
$\big(\it{\Omega}^1_{X/Y}\big)_{x}$ form a basis of this
module. Taking $X$ small enough, one may assume that the $g_{i}$
come from sections of~$\cal{O}_{X}$, and thus define a
$Y$-morphism $g\colon X\to Y'=Y[t_{1},\dots,t_{n}]$. Using
the hypothesis and lemma~\ref{II.4.2}, one sees easily that
the corresponding homomorphism
$g^*\big(\it{\Omega}^1_{Y'/S}\big)\to\it{\Omega}^1_{X/S}$ is
bijective at~$x$, which reduces us to proving the

\begin{corollary}
\label{II.4.6}
Let $f\colon X\to Y$ be a morphism of smooth $S$-preschemes.
For $f$ to be \'etale at $x\in X$, it is necessary and sufficient
that $f^*\big(\it{\Omega}^1_{Y/S}\big)\to\it{\Omega}^1_{X/S}$ be
an isomorphism at~$x$.
\end{corollary}

One knows that this is necessary by~\ref{II.4.1}, and this condition
implies that $f$ is unramified at~$x$ by the same lemma. By virtue
of~\ref{II.2.2}, one is reduced to the case where
$S=\Spec(k)$. Since $Y$ is smooth over~$k$, it is regular, hence a
fortiori normal, and by virtue of (I~\SourceRef{I.9.5}{9.5}~(ii)) one is reduced
to proving that $\cal{O}_{y}\to\cal{O}_{x}$ is injective, or equivalently
that $\cal{O}_{y}$ and $\cal{O}_{x}$ have the same dimension. Now these
dimensions are respectively the ranks of~$\it{\Omega}^1_{Y/k}$ and
$\it{\Omega}^1_{X/k}$ at~$y$ \resp $x$, hence equal by virtue of
the hypothesis.
\begin{remarks}
\label{II.4.7}
$X$ and $Y$ being assumed smooth over~$S$, the
criterion~\ref{II.4.5}~(ii) of smoothness of $f\colon X\to Y$ can
also be stated by saying that for every $x\in X$, the
\emph{tangent} map (relative to the base~$S$) of~$f$ at~$x$,
\ie the transpose of the homomorphism of the $\kres(x)$ finite-dimensional vector spaces,
restricted fibers of
$f^*\big(\it{\Omega}^1_{Y/S}\big)$ and $\it{\Omega}^1_{X/S}$ at
$x$, is \emph{surjective}. This is a quite
familiar hypothesis in particular among people working with
analytic spaces. The non-singularity hypothesis which they
ordinarily make (which means that $X$ and $Y$ are ``smooth over~$\CC$'', \cf \No\ref{II.5}) seems due only to the fear which
singular points of algebraic varieties or analytic spaces still inspire
in many geometers.
\end{remarks}

Let us point out the following particular case of~\ref{II.4.6}:
\begin{corollary}
\label{II.4.8}
Let $X$ be an $S$-prescheme, $g\colon X\to
S[t_{1},\dots,t_{n}]$ an $S$-morphism, defined by the sections
$g_{i}$ ($1\le i\le n$) of~$\cal{O}_{X}$, $x$ a point of~$X$ such that
$X$ is smooth over~$S$ at~$x$. For $g$ to be \'etale at~$x$, it
is necessary and sufficient that the $\rd g_{i}$
% original p. 39
($1\le i\le n$) form a basis of~$\it{\Omega}^1_{X/S}$ at~$x$ (or,
which comes to the same thing, that their images in
$\it{\Omega}^1_{X/S}(x) =
\big(\it{\Omega}^1_{X/S}\big)_{x}\otimes_{\cal{O}_{x}}\kres(x)$ form
a basis of this vector space over~$\kres(x)$).
\end{corollary}

Let $X$ be a prescheme, $Y$ a closed subprescheme
of~$X$ defined by a coherent sheaf $\cal{J}$
of ideals. Thus $\cal{J}/\cal{J}^2$ may be regarded
as a coherent sheaf on~$Y$ (the \emph{conormal sheaf}
of~$Y$ in~$X$). If now $X$ is an $S$-prescheme, one has
a canonical exact sequence of quasi-coherent sheaves on~$Y$
\begin{equation}
\label{eq:II.4.3}
\cal{J}/\cal{J}^2\lto{\rd}
\it{\Omega}^1_{X/S}\otimes_{\cal{O}_{X}}\cal{O}_{Y}\to
\it{\Omega}^1_{Y/S}\to 0
\end{equation}
whose right-hand part is none other than~\eqref{eq:II.4.2} (with the
role of~$X$ and~$Y$ interchanged, taking into account that
$\it{\Omega}^1_{Y/X}=0$), while the homomorphism
$\cal{J}/\cal{J}^2\to\it{\Omega}^1_{X/S}\otimes_{\cal{O}_{X}}\cal{O}_{Y}$
is deduced from the homomorphism (in general not
linear) $g\to\rd g$ by passing to
quotients. The exactness of~\eqref{eq:II.4.3} is classical and
moreover trivial, and is interpreted in the affine case by the
following exact sequence (corresponding to a surjective homomorphism
$B\to C$ of $A$-algebras, with kernel~$J$):
\begin{equation*}
\label{eq:II.4.3bis}
\tag{4.3\textrm{bis}}
J/J^2\to\Omega^1_{B/A}\otimes_{B}C\to\Omega^1_{C/A}\to 0\qquad
(C=B/J)
\end{equation*}
(an exact sequence which had already been used
implicitly in the proof of (I~\SourceRef{I.7.2}{7.2})!).

\begin{proposition}
\label{II.4.9}
Let $X$ be an $S$-prescheme, $Y$ a closed subprescheme
of~$X$ defined by a coherent sheaf $\cal{J}$
of ideals on~$X$, $x$ a point of~$X$, $g_i$ ($1 \leq i \leq n$)
sections of~$\cal{O}_X$, defining an $S$-morphism
\[
g \colon X \to S[t_1,\dots,t_n] = X'
\]
finally let $p$ be an integer, $0 \leq p \leq n$. Suppose $X$ is
\emph{smooth over~$S$ at $x$}. The following conditions are equivalent:
\begin{enumerate}
\item[(i)] There exists an open neighborhood $X_1$ of~$x$ in $X$ such that
$g|X_1$ is \emph{\'etale} and $Y_1=Y \cap X_1$ (trace of~$Y$
on~$X_1$) is the \emph{inverse image} of the closed subprescheme
$Y'=S[t_{p+1},\dots,t_n]$
of
$X'=S[t_1,\dots,t_n]$ (\ie the $g_i$ ($1 \leq i \leq p$) generate
$\cal{J}|X_1$):
\[
\xymatrix{ Y_1 \ar[d] \ar[r] & X_1 \ar[d]^{\hbox{\'etale}} \\
Y'=S[t_{p+1},\dots, t_n] \ar[r] & X'=S[t_1,\dots,t_n]}
\]
\item[(ii)]
% original p. 40
$Y$ is \emph{smooth over~$S$ at $x$}, the $g_i$ ($1 \leq i \leq p$)
define \emph{elements of} $\cal{J}_x$, the $d
g_i(x)$ ($1 \leq i \leq n$) form a \emph{basis of}
$\it{\Omega}^1_{X/S}(x)$ over~$\kres(x)$, the $dg'_i(x)$ ($p+1 \leq i \leq
n$) form a \emph{basis of} $\it{\Omega}^1_{Y/S}(x)$ over~$\kres(x)$
(where the $g'_i$ denote the restrictions of the $g_i$ to~$Y$;
the differentials are taken with respect to~$S$).
\item[(iii)] The $g_i$ ($1 \leq i \leq p$) define a
\emph{system of generators} of~$\cal{J}_x$, and the
$dg_i(x)$ ($1 \leq i \leq n$) form a \emph{basis of}
$\it{\Omega}^1_{X/S}(x)$ over~$\kres(x)$.
\item[(iv)] $Y$ is \emph{smooth over~$S$} at $x$, the $g_i$
($1 \leq i \leq p$)
form a \emph{minimal system of generators of}
$\cal{J}_x$, the $dg'_i(x)$ ($p+1 \leq i \leq n$) form a
\emph{basis of} $\it{\Omega}^1_{Y/S}(x)$ over~$\kres(x)$.
\end{enumerate}

Moreover, under these conditions, $\cal{J}/\cal{J}^2$ is a free Module
on~$Y$ at $x$, admitting as \emph{basis at $x$} the classes of the
$g_i$ ($1 \leq i \leq p$), \emph{and the canonical homomorphism}
$\cal{J}/\cal{J}^2 \to \it{\Omega}^1_{X/S} \otimes \cal{O}_Y$ is
\emph{universally injective at $x$}.
\end{proposition}

\begin{remarkstar}
This implies that $p$ is well determined by the other
conditions, either as the \emph{rank} of the free Module
$\cal{J}/\cal{J}^2$ on~$Y$ at $x$, or again as the \emph{minimum
number of generators} of~$\cal{J}_x$ on~$X$, or finally by the fact
that the relative dimension of~$Y$ rel.\ $S$ at $x$ is $n-p$.
\end{remarkstar}

\subsubsection*{Proof}
Suppose first that (i) is satisfied. Then by (I~\SourceRef{I.4.6}{4.6}
(iii)) $Y_1$ is \'etale over~$Y'$, hence by definition it is
smooth over~$S$ at $x$ (of relative dimension $n-p$), hence the same
is true of~$Y$. It then follows from \eqref{II.4.8} that the $dg_i$
($1 \leq i \leq n$) form a basis of~$\it{\Omega}^1_{X/S}$ at $x$,
and that the $dg'_i$ ($p+1 \leq i \leq n$) a basis of
$\it{\Omega}^1_{Y/S}$ at $x$, whence it follows by the exact
sequence \eqref{eq:II.4.3} that the $g_i$ ($1 \leq i \leq p$) are
linearly independent in $\cal{J}/\cal{J}^2$
(regarded as a Module on~$Y$) at $x$; since the $g_i$ ($1
\leq i \leq p$) generate $\cal{J}_x$, it follows that the $g_i \mod
\cal{J}_x^2$ form a \emph{basis of} $\cal{J}/\cal{J}^2$ at
$x$. This implies on the one hand that the $g_i$ ($1 \leq i \leq p$)
form a \emph{minimal} system of generators of
$\cal{J}_x$, on the other hand that the homomorphism $\cal{J}/\cal{J}^2 \to
\it{\Omega}^1_{X/S} \otimes \cal{O}_Y$ of \eqref{eq:II.4.3} is
universally injective at $x$ (for it maps a basis of a free Module at $x$
onto a part of a basis of a free Module at $x$ --- N.B.\ these are~$Y$-Modules). This proves that (i) implies (ii),
(iii), (iv), as well as the last assertions of
proposition~\ref{II.4.9}.

(iii) implies (i) by virtue of corollary~\ref{II.4.8}.

(ii)
% original p. 41
implies (i). Indeed, the first hypothesis in (ii)
means that (replacing if necessary $X$ by an open neighborhood of
$x$ in $X$) $g$ induces a morphism $h\colon Y \to
Y'$. By~\ref{II.4.8}, the other two hypotheses of (ii)
mean that $g$ is \'etale at $x$, and $h$ \'etale at $x$. Let
then $Y''$ be the inverse image of~$Y'$ by $g$. Thus $Y$ is a
closed subprescheme of~$Y''$, which is \'etale over~$Y'$
at $x$ by (I~\SourceRef{I.4.6}{4.6}~(iii)) since $g$ is \'etale at
$x$. Hence the immersion morphism $Y \to Y''$ is itself
\'etale (I~\SourceRef{I.4.8}{4.8}) hence an open immersion (I~\SourceRef{I.5.8}{5.8} or
I~\SourceRef{I.5.2}{5.2}), hence replacing still $X$ by a suitable open
neighborhood $X_1$ of~$x$, one obtains (i).

What precedes establishes the equivalence of conditions
(i), (ii), (iii), and the fact that they imply (iv), it remains to
prove that (iv)$\To$(ii), which is immediate (taking into
account that $\it{\Omega}^1_{X/S}$ is free on~$X$ at $x$) once
one knows that the fact that $Y$ is smooth over~$S$ at $x$ implies that
$\cal{J}/\cal{J}^2$ is free on~$Y$ at $x$, and the homomorphism
$\cal{J}/\cal{J}^2 \to \it{\Omega}^1_{X/S}\otimes \cal{O}_Y$
universally injective at $x$. This last point is included in the
\begin{theorem}
\label{II.4.10}
Let $X$ be a smooth $S$-prescheme, $Y$ a
closed subprescheme of~$X$ defined by a coherent sheaf
$\cal{J}$ of ideals on~$X$, $x$ a point of~$X$. The
following conditions are equivalent:
\begin{enumerate}
\item[(i)] $Y$ is \emph{smooth over~$S$ at $x$}.
\item[(ii)] There exists an open neighborhood $X_1$ of~$x$ in $X$ and an
\emph{\'etale} $S$-morphism
\[
g\colon X_1 \to X'=S[t_1,\dots,t_n]
\]
such that $Y_1=Y \cap X_1$ (trace of~$Y$ on~$X_1$) is the
subprescheme of~$X_1$ \emph{inverse image} by $g$ of the
closed subprescheme $Y'=S[t_{p+1},\dots,t_n]$ of
$X'=S[t_1,\dots,t_n]$, for a suitable $p$.
\item[(iii)] There exist \emph{generators $g_i$ ($1 \leq i
\leq p$) of~$\cal{J}_x$}, such that the $dg_i$ form part of a
basis of~$\it{\Omega}^1_{X/S}$ at $x$ (or, what amounts to the same,
such that the $dg_i(x)$ in $\it{\Omega}^1_{X/S}$ are
linearly independent over~$\kres(x)$).
\item[(iv)] The sheaf $\cal{J}/\cal{J}^2$ is free on~$Y$ at $x$,
and the canonical homomorphism
\[
d\colon \cal{J}/\cal{J}^2 \to \it{\Omega}^1_{X/S}\otimes \cal{O}_Y
\]
is universally injective at~$x$; or again: the sequence
of canonical homomorphisms
\[
0 \to \cal{J}/\cal{J}^2 \to \it{\Omega}^1_{X/S}\otimes \cal{O}_Y \to
\it{\Omega}^1_{Y/S} \to 0
\]
is exact at $x$, and $\it{\Omega}^1_{Y/S}$ is locally free
at~$x$.
\end{enumerate}
\end{theorem}

\subsubsection*{Proof}
% original p. 42
One already knows that (ii) implies (i), (iii), (iv) by what
precedes. Let us prove that (i)$\To$(ii) (which
will at the same time complete the proof
of~\ref{II.4.9}).
By theorem~\ref{II.4.3}~(ii), the last two
terms in the exact sequence \eqref{eq:II.4.3} are free Modules
on~$Y$. Thus, since the images in
$\it{\Omega}^1_{X/S}\otimes_{\cal{O}_X} \cal{O}_Y$ of the $dg$ ($g \in
\cal{O}_X$) generate this module at $x$, hence their images in
$\it{\Omega}^1_{Y/S}$ generate the latter at~$x$, one can find
$g_i$ ($p+1 \leq i \leq n$) in $\cal{O}_X$ such that the $dg'_i$
form a basis of~$\it{\Omega}^1_{Y/S}$, then (by virtue of
the exactness of \eqref{eq:II.4.3}) complete the system of
$dg_i$ ($p+1 \leq i \leq n$) to a basis of the middle term by
elements of the form $dg_i$ ($1 \leq i \leq n$) where the
$g_i$ ($1 \leq i \leq p$) \emph{are in $\cal{J}_x$.} The $g_i$
come from sections of~$\cal{O}_X$ on a neighborhood of~$x$ in
$X$, which one may suppose equal to $X$. One is then under the
conditions of~\ref{II.4.8}~(ii), and one has established that this implies
condition~\ref{II.4.8}~(i), whence~\ref{II.4.10}~(ii).

The implication (iii)$\To$(ii) in~\ref{II.4.10} follows
at once from the implication (iii)$\To$(i)
in~\ref{II.4.8}. Hence
(i), (ii), (iii)
are equivalent, and
imply (iv). Finally, it is trivial that (iv) implies (iii), taking
into account that $g_i \in \cal{J}_x$ which form a basis of~$\cal{J}_x
\mod \cal{J}_x^2$ generate $\cal{J}_x$ (Nakayama).

Moreover, the preceding proof shows the following:
\begin{corollary}
\label{II.4.11}
Let $X$ be an $S$-prescheme, $Y$ a closed subprescheme
defined by a coherent sheaf $\cal{J}$
of ideals on~$X$, $x$ a point of~$Y$. Suppose \emph{$X$ and $Y$
are smooth over~$S$ at $x$}. Let $g_i$ be sections of~$\cal{J}$ ($1
\leq i \leq p$). The following conditions are equivalent:
\begin{enumerate}
\item[(i)] The $g_i$ \emph{generate} $\cal{J}_x$ and the $dg_i(x)$
are \emph{linearly independent} in
$\it{\Omega}^1_{X/S}(x)$
over~$\kres(x)$.
\item[(ii)] The $g_i \mod \cal{J}^2$ form a basis of
$\cal{J}/\cal{J}^2$ at $x$.
\item[(iii)] The $g_i$ form a minimal system of
generators of~$\cal{J}_x$.
\item[(iv)] One can find other sections $g_i$ ($p+1 \leq i \leq
n$) of~$\cal{O}_X$ on a neighborhood $X_1$ \emph{of} $X$,
defining with the preceding ones an
\emph{\'etale} morphism $X_1 \to X'=S[t_1,\dots,t_n]$ such that $Y_1=Y \cap
X_1$ is the \emph{inverse image} by $g$ of the closed subprescheme
$Y'=S[t_{p+1},\dots,t_n]$ of~$X'=S[t_{1},\dots,t_n]$.
\end{enumerate}
\end{corollary}

% original p. 43
In particular:

\begin{corollary}
\label{II.4.12}
Let $X$ be an $S$-prescheme, $F$ a section of~$\cal{O}_X$,
$Y$ the subprescheme of zeros of~$F$ (defined by
the coherent Ideal $F\cdot\cal{O}_X$), $x$ a point of~$Y$.
Suppose $X$ is simple over~$S$ at $x$. For $Y$ to be smooth over~$S$ at
$x$, it is necessary and sufficient that either $F$ be zero in a neighborhood of
$x$, or that $dF(x) \neq 0$ (where $dF(x)$ denotes the image of
$dF$ in the vector space $\it{\Omega}^1_{X/S}(x)$ over~$\kres(x)$).
\end{corollary}

This is sufficient by virtue of~\ref{II.4.10} criterion~(iii). This is
necessary, for since $\cal{J}$ is generated by one
element, one must first have that $\cal{J}/\cal{J}^2$ at the point $x$
be free of rank $\leq 1$. If this rank is $0$,
\ie $\cal{J}/\cal{J}^2=0$ at $x$, it follows that $\cal{J}=0$ at $x$
by Nakayama, \ie $F$ is zero in a neighborhood of~$x$. If this rank is
$1$, then $F$ forms a minimal system of generators of
$\cal{J}$ at $x$, and one concludes by \eqref{II.4.11}, equivalence of
(i) and~(iii)).

\begin{corollary}
\label{II.4.13}
Let $Y$ be an $S$-prescheme locally of finite type, $S'$ a
\emph{flat} $S$-prescheme, $Y'=Y \times_S S'$, $x'$ a point
of~$Y'$ and $x$ its canonical image in~$Y$. For $Y$ to be smooth
over~$S$ at~$X$, it is necessary and sufficient that $Y'$ be smooth over~$S'$ at~$x'$. In particular, if $S'\to S$ is flat and surjective, $Y$ is
smooth over~$S$
if and only if $Y'$ is smooth over~$S'$.
\end{corollary}

Only sufficiency needs to be proved (necessity was
seen in~\ref{II.1.3}~(ii)). One may suppose (replacing $Y$ by a suitable neighborhood of~$x$, $Y'$ by the inverse image
of the latter) that $Y$ is affine of finite type over~$S$ affine, hence
$Y$ is isomorphic to a closed subprescheme of a
scheme $S[t_1,\dots,t_n]$. Consequently, $Y'$ is identified with a
closed subprescheme of~$X'=X \times_S S'$. Since $X$ is
smooth over~$S$, hence $X'$ smooth over~$S'$, one can apply the
smoothness criteria~\ref{II.4.10}. Here, criterion (iv)
gives the result at once.

\begin{remarks}
\label{II.4.14}
Criterion (iii) of Theorem~\ref{II.4.10} deserves to be called
the \emph{Jacobian criterion of smoothness}.
It allows one to recognize, theoretically, whether a given
$S$-prescheme $Y$ is smooth over~$S$ at a point $x$
of~$Y$, since there always exists a neighborhood of~$Y$ isomorphic to
a subprescheme of a smooth $S$-prescheme $X$, for
example $X=S[t_1,\dots,t_n]$. It is moreover for
$X=S[t_1,\dots,t_n]$, $S=\Spec(A)$, that one usually states the
Jacobian criterion
% original p. 44
(of course, in the classical case envisaged by Zariski, $A$ was
a field). One leaves to the reader to give the statement relative
to the datum of an ideal $J$ of~$A[t_1,\dots,t_n]$ and of a
prime ideal containing it, to which one is thus led. Let us note
that it indeed seems at the present time (and especially since Nagata
has succeeded in generalizing by
non-differential methods Zariski's theorem stating that
the set of regular points of an algebraic scheme is
open) that the Jacobian criterion is scarcely of interest
except in the form in which we give it here (\ie using
exclusively \emph{relative} differentials and not
\emph{absolute} differentials, \ie relative to the ring of
absolute constants $\ZZ$). As so often, the consideration
of differentials is more convenient here than that of
derivations. Let us note finally that if $Y$ is smooth over~$S$ at $x$, of
relative dimension $n$, then there exists an open neighborhood of~$x$
in~$Y$ isomorphic to a subprescheme of
$X=S[t_1,\dots,t_n]$ with $n=m+1$, as follows from the
definition and
from~(I~\SourceRef{I.7.6}{7.6}).

Let $A$ be a noetherian ring, $x_i$ ($1 \leq i \leq n$)
elements of~$A$, $J$ the ideal generated by the
$x_i$. One says that the $x_i$ form a \emph{regular
system of generators}
of~$J$ if the canonical surjective homomorphism
\[
(A/J)[t_1,\dots,t_n] \to \gr^J(A)
\]
defined by the $x_i$ (where the second member denotes
the graded ring associated to $A$ filtered by the
powers of~$J$) is an \emph{isomorphism}. This condition
also means that
\begin{enumerate}
\item[(i)] The canonical surjective homomorphism
\[
S_{A/J} (J/J^2) \to \gr^J(A)
\]
(where the first member denotes the symmetric algebra of the
$A/J$-module $J/J^2$) is an isomorphism, and
\item[(ii)] $J/J^2$ is free and admits as basis the classes of the $x_i
\mod J^2$.
\end{enumerate}
In this form, one sees that if $J\neq A$, the $x_i$ form a
\emph{minimal system of generators} of~$J$, and that
\emph{every other minimal system of generators} of~$J$
\emph{is a regular system of generators}
(N.B.\ ``minimal'' is taken in the strict sense: minimum number
of elements, which is equivalent to the sense: minimal for
inclusion, only if $A$ is local); on the other hand, if $J=A$, every
system of generators of~$J$ is regular.

The
% original p. 45
condition of regularity of a system of
generators of an ideal is stable under localization with respect to an
arbitrary multiplicatively closed set, and on the other hand one sees
at once that for $(x_i)$ to be a minimal system of
generators of~$J$, it already suffices that for every
\emph{maximal ideal $\goth{m}$ containing $J$}, the $x_i$ define a
regular system of generators of~$JA_{\goth{m}}$ in
$A_{\goth{m}}$. This therefore reduces us to the case where $A$ is a
local ring with maximal ideal $\goth{m}$, and where the $x_i$ are
in $\goth{m}$. \emph{Then the $x_i$ form a
regular system of generators of~$J$ if and only if they
form an $A$-sequence in the sense of Serre}\footnote{We shall
now rather say ``$A$-regular sequence'', \cf EGA
$0_{\textup{IV}}$~15.1.7 and 15.1.11.}, \ie if for every $i$ such that $1
\leq i \leq n$, $x_i$ is not a divisor of~$0$ in
$A/(x_1,\dots,x_{i-1})A$.

Finally, in the case where $A$ is an algebra over a ring $B$,
and where $A/J$ is isomorphic as a $B$-algebra to $B$ (so
that $J$ is the kernel of a homomorphism of $B$-algebras
$A\to B$), then the $x_i$ form a regular system of
generators of~$J$ if and only if the canonical
homomorphism
\[
B[[t_1,\dots,t_n]]\to \hat{A}
\]
defined by the $x_i$ (where the second member denotes the
separated completion $\varprojlim A/J^{n+1}$ of~$A$ for the
topology defined by the powers of~$J$) is an
\emph{isomorphism} (it is in any case \emph{surjective}).

All these facts are well known, and doubtless appear in Serre's
course on commutative algebra written up by Gabriel, up
to minor details. (Where one finds $N$ other
characterizations of $A$-sequences, in the case where $A$ is a
local ring).

Let $J$ be an ideal in a noetherian ring $A$. We shall say that
$J$ is a \emph{regular ideal}
if for every prime ideal $\goth{p}$ of~$A$, $JA_{\goth{p}}$ admits
a regular system of generators. It is
obviously sufficient to verify this for $\goth{p} \supset J$, and one may
moreover restrict oneself to maximal $\goth{p}$. More generally,
let $\cal{J}$ be an ideal on a locally
noetherian prescheme $X$, one says that $\cal{J}$ is a \emph{regular
ideal} if for every $x\in X$, $\cal{J}_x$ is an ideal of
$\cal{O}_x$ which admits a regular system of
generators. This is equivalent to the conjunction of the two
following conditions:
\begin{enumerate}
\item[(a)] The canonical surjective homomorphism
\[
S_{\cal{O}_X/\cal{J}}(\cal{J}/\cal{J}^2) \to \gr^{\cal{J}}(\cal{O}_X)
\]
is an isomorphism and
\item[(b)] The sheaf of~$\cal{O}_X/\cal{J}$-Modules
$\cal{J}/\cal{J}^2$ is locally free.
\end{enumerate}

One
% original p. 46
then also says that the subprescheme $Y$ of~$X$ defined
by $\cal{J}$ (hence such that $\cal{O}_Y$ extended by $0$ is
isomorphic to $\cal{O}_X/\cal{J}$) is \emph{regularly
immersed}
in~$X$, and one defines in the same way (in an evident
manner) the notion of morphism of \emph{regular
immersion},
(\resp \emph{regular at a point $x$}),
immersion morphism $Y\to X$ identifying $Y$ (\resp a suitable neighborhood of
$x$) with a closed subprescheme regularly
immersed in an open subset of~$X$. (One must not say:
regular subprescheme, for that would mean that the
local rings of~$Y$ are regular). Finally, sections $x_i$
of~$\cal{J}$ are called a \emph{regular system of
generators} if for every $x\in X$, the corresponding
elements of~$\cal{O}_x$ form a regular system of
generators of~$\cal{J}_x$ (terminology compatible with that
introduced for generators of an ideal of a
ring). This also means that the canonical surjective homomorphism
\[
\cal{O}_Y[t_1,\dots,t_n] \to \gr^{\cal{J}} (\cal{O}_X)
\]
defined by the $x_i$ is an isomorphism. If one knows in advance
that the ideal $\cal{J}$ is regular, this also means,
simply, that at every point $x$ \emph{of~$Y$}, the $x_i$
define a \emph{basis} of~$\cal{J}/\cal{J}^2$ over~$\cal{O}_{Y,x}$. (N.B.\ this condition is empty if $Y$ is
empty). Thus, for $\cal{J}$ to admit a regular
system of generators, it is necessary and sufficient that $\cal{J}$ be
regular, and that the $\cal{O}_Y$-Module $\cal{J} /\cal{J}^2$ be
globally free (and not only locally free), \ie that
the canonical homomorphism $S_{\cal{O}_Y}(\cal{J}/\cal{J}^2) \to
\gr^{\cal{J}} (\cal{O}_X)$ be surjective, and that the
$\cal{O}_Y$-Module $\cal{J}/\cal{J}^2$ be globally free.

An \emph{augmented ring is said to be regular}
if the ideal of the augmentation is regular. Thus, if $A$ is
a local ring, regarded as augmented in its residue
field $k$, then $A$ is a regular local ring if and
only if it is a regular augmented ring.

(To tell the truth, it seems that it was useless to begin by
doing the preliminary sorites for rings; there is an
advantage in beginning with sheaves at once. If
one wants something in the noetherian case, it is the
definition adopted here---a priori less strict than that by
Serre's $A$-sequences---which seems preferable for the
needs of differential calculus. Of course, to do things properly, one
should also develop at least part of the theory
of smooth morphisms in the non-noetherian setting\footnote{As
is said in the foreword, this has now been done, \cf EGA
IV 17, 18.}, probably
% original p. 47
starting from the Jacobian criterion, so as to obtain if
possible all the essential formal properties of
smooth morphisms and of \'etale morphisms \ie smooth and
quasi-finite; the converses alone appealing to
noetherian hypotheses).
\end{remarks}

After these long terminological preliminaries, a little
theorem:

\begin{theorem}
\label{II.4.15}
Let $X$ be an $S$-prescheme locally of finite type, $Y$ a
closed subprescheme of~$X$ defined by a coherent sheaf
$\cal{J}$ of ideals on~$X$, $x$ a point of~$X$. One
now assumes \emph{$Y$ is smooth over~$S$ at $x$} (and nothing on~$X$). Then the following conditions are equivalent:
\begin{enumerate}
\item[(i)] $X$ is smooth over~$S$ at $x$
\item[(ii)] The immersion $i\colon Y \to X$ is regular at
$x$, \ie $\cal{J}_x$ is a regular ideal of~$\cal{O}_x$.
\end{enumerate}
\end{theorem}

\begin{corollary}
\label{II.4.16}
Suppose $Y$ is \emph{smooth} over~$S$. For $X$ to be smooth over~$S$
in a neighborhood of~$Y$ (\ie at the points of~$Y$) it is necessary and sufficient
that $Y$ be regularly immersed in $X$, \ie that
the immersion $i \colon Y \to X$ be regular.
\end{corollary}

\subsubsection*{Proof}
(i) implies (ii). One applies \ref{II.4.10} criterion~(ii), since
$g\colon X_1 \to X$ is \emph{flat}, to show that the inverse image
under $g$ of the subprescheme $Y'$ of~$X'$ is
regularly immersed, one is reduced to proving that
$Y'=S[t_{p+1},\dots, t_n]$ is regularly immersed in
$S[t_1,\dots,t_n]$, which is trivial (the $t_i$ ($1 \leq i \leq
p$) form a regular system of generators of
the Ideal defining $Y'$ in $X'$).

(ii) implies (i). Let $g_i$ ($1 \leq i \leq p$) be a regular
system of generators of~$\cal{J}_x$ and let $g_i$
($p+1 \leq i \leq n$) be elements of~$\cal{O}_{X,x}$ such that
their images $g'_i$ in $\cal{O}_{Y,x}$ define an
\emph{\'etale} morphism
\[
Y_1 \to Y'=S[t_{p+1},\dots, t_n]
\]
of a neighborhoods $Y_1$ of~$Y$ in $Y'$. The $g_i$ ($1 \leq i \leq n$)
come from sections (of the same name) of~$\cal{O}_X$ on a
neighborhood $X_1$ of~$x$, and one may suppose $X_1=X$, $Y_1=Y$. One
thus obtains a morphism
\[
g\colon X \to X'=S[t_1,\dots,t_n]
\]
and everything reduces to showing that this morphism is \emph{\'etale}
at~$x$. Taking $X_1$ small enough,
% original p. 48
one may suppose that the $g_i$ ($1 \leq i \leq p$) form a
regular system of generators of~$\cal{J}$ on all of
$X$. In particular, they generate $\cal{J}$, hence the
subprescheme $Y$ of~$X$ is identified with the inverse image
under $g$ of the subprescheme $Y'$ of~$X'$. Let $x'=g(x)$, then
the fiber of~$X' \to X$ at $x'$ is
identical to the fiber of~$Y \to Y'$ at $x$, hence is \'etale over~$\kres(x')$, hence $g$ is \emph{unramified} at $x$, it remains to
prove that $g$ is \emph{flat} at $x$. Now the associated graded
of $\cal{O}_{X',x'}$ filtered by the powers of
$\cal{J}'_{x}$ is \emph{free} over~$\cal{O}_{Y',x'}$ in all
degrees, on the other hand the associated graded of
$\cal{O}_{X,x}$ filtered by the powers of
$\cal{J}_x=\cal{J}'_x\cal{O}_{X,x}$ is isomorphic (under
the canonical homomorphism) to the tensor product of the former
with $\cal{O}_{Y,x}$ (since both rings are polynomial
rings in $n-p$ indeterminates, with ring of
constants $\cal{O}_{Y',x'}$, \resp $\cal{O}_{Y,x}$), finally over~$\cal{O}_{X',x'}/\cal{J}'_{x'}= \cal{O}_{Y',x'}$
$\cal{O}_{X,x}/\cal{J}_x =\cal{O}_{Y,x}$ is flat.

By a general flatness criterion (valid for
a local homomorphism of noetherian local rings $A' \to A$,
$A'$ being equipped with an ideal $J' \neq A'$ such that the associated
graded is free over~$A'/J'$ in every dimension) it follows that
$X$ is flat over~$X'$ at $x$, qed.

\begin{corollary}
\label{II.4.17}
Let $X$ be a prescheme locally of finite type over~$Y$, $i$
a section of~$X$ over~$Y$, $y$ a point of~$Y$, $x=i(y)$, $\cal{J}$
the sheaf of ideals on~$X$ defined by the
subprescheme $i(Y)$ (which we assume closed to
simplify the statement, a condition satisfied if $X$ is a
scheme).

The following conditions are equivalent:
\begin{enumerate}
\item[(i)] $X$ is smooth over~$Y$ at $x$
\item[(ii)] $i$ is a regular immersion at $y$
\item[(iii)] The completed $\cal{O}_y$-algebra of
$\cal{O}_x$ for the topology defined by the powers of
$\cal{J}_x$ is isomorphic to a formal power
series algebra $\cal{O}_y[[ t_1,\dots,t_n]]$.
\item[(iii~bis)] There exists an open neighborhood $U$ of~$y$ such that
the sheaf of algebras\linebreak
$\varprojlim i^*(\cal{O}_X/\cal{J}^{n+1})$
over~$\cal{O}_Y$ is isomorphic to a sheaf of the form
$\cal{O}_Y[[ t_1,\dots, t_n]]$ over~$U$.
\item[(iv)] There exists an open neighborhood $U$ of~$y$, and an
open neighborhood $V$ of~$x$, and finally a $Y$-morphism $g \colon V \to
U[t_1,\dots, t_n]$, such that $g$ is \'etale, that
$i$ induces a section of~$V$ over~$U$, transformed by $g$ into the
zero section of~$U[t_1,\dots, t_n]$ over~$U$.
\end{enumerate}
\end{corollary}

The equivalence
% original p. 49
of (i) and (ii) is a special case of
theorem~\ref{II.4.15}, setting $Y=S$, the equivalence of
(ii) and (iii) (and morally of (ii) and (iii~bis)) has been
pointed out with the ``recalls''. as for the equivalence of (i)
and (iv), it is easily deduced from
theorem~\ref{II.4.10} (equivalence of conditions (i) and
(ii) of the latter).

\begin{corollary}
\label{II.4.18}
Let $X$ be a prescheme smooth over~$S$. Then the
diagonal morphism
\[
\Delta_{X/S}\colon X\to X\times_S X
\]
is a
\emph{regular immersion},
or as one also says, $X$ is ``\emph{differentially
smooth}'' over~$S$.
\end{corollary}

Indeed, this is a special case of corollary~\ref{II.4.16}, since
$X$ and $X\times_S X$ are both smooth over~$S$.

% The corrected source repeats 4.18. Preserve it with a distinct PDF anchor.
\setcounter{theorem}{17}
\begingroup
\renewcommand{\theHtheorem}{\theHsection.\arabic{theorem}.bis}
\begin{remarks}
\label{rem:II.4.18}
Let us recall (I~\SourceRef{I.1}{1}) that if $X$ is a prescheme over
$S$, one introduces the quasi-coherent sheaves
of algebras
$\cal{P}_{X/S}^n=\cal{O}_{X\times_S X}/\cal{I}_X^{n+1}$ on~$X$,
(where $\cal{I}_X$ denotes the sheaf
of ideals
which defines the diagonal in $X\times_S X$), considered as
a sheaf of~$\cal{O}_X$-algebras thanks to the first
projection $\pr_1\colon X\times_S X\to X$. The
$\cal{P}_{X/S}^n$ form a projective system of Algebras over~$X$, whose projective limit is denoted $\cal{P}_{X/S}^{\infty}$
and is none other than the structure sheaf of the formal completion
of~$X\times_S X$ along the diagonal (now assuming $X$ is
locally of finite type over~$S$, hence the $\cal{P}_{X/S}^n$
are coherent). To say that $X$ is differentially smooth over~$S$
(\ie that the diagonal morphism $\Delta_{X/S}$ is a regular
immersion) also means that $\cal{P}_{X/S}^{\infty}$ is
regular, as a sheaf of algebras augmented to
$\cal{O}_X$, \ie that $\it{\Omega}_{X/S}^1$ is locally free and
the canonical surjective homomorphism
\[
S_{\cal{O}_X}(\it{\Omega}_{X/S}^1)\to
\mathrm{gr}_*(\cal{P}_{X/S}^{\infty})
\]
is an isomorphism, or finally that every point of~$X$
has an open neighborhood on which the sheaf of
augmented algebras $\cal{P}_{X/S}^{\infty}$ is isomorphic to a
sheaf $\cal{O}_X[[t_1,\dots,t_n]]$.

Let $s$ be a section of~$X$ over~$S$, $\cal{J}$ the sheaf
of ideals on~$X$ which it defines (assuming
% original p. 50
for simplicity that $s(S)$ is closed), one then has
canonical isomorphisms of augmented~$\cal{O}_X$-algebras:
\begin{equation}
\label{eq:II.4.4}
s^*(\cal{P}_{X/S}^n)=\cal{O}_X/\cal{J}^{n+1}\quoi,\quad
s^*(\cal{P}_{X/S}^{\infty})= \varprojlim_{n}
\cal{O}_X/\cal{J}^{n+1}
\end{equation}
These isomorphisms are functorial in an evident sense under
change of base, and (taking this fact into account) give back a
characterization of the sheaves of algebras $\cal{P}_{X/S}^n$ over~$S$. If for example $S=\Spec(k)$, $k$ a field, then the datum
of a section~$s$ of~$X$ over~$S$ is equivalent to the datum of a
point $x$ of~$X$ rational over~$k$, and the
preceding formulas mean that one has an isomorphism of
$k$-algebras
\begin{equation}
\label{eq:II.4.5}
\cal{P}_{X/S}^n(x)=\cal{O}_x/\goth{m}_x^{n+1}
\end{equation}
which justifies the name: ``\emph{sheaf of principal parts
of order~$n$ on~$X$ rel.\ to~$S$}''
given to~$\cal{P}_{X/S}^n$. One sees moreover from
\eqref{eq:II.4.4} that \emph{if $X$ is differentially smooth
over~$S$ at every point of~$s(S)$, then $X$ is smooth over~$S$ at every
point of~$s(S)$}, (corollary~\ref{II.4.17}) \emph{the converse
being equally true} (corollary~\ref{II.4.18}). Taking
\ref{II.4.13} into account, one easily concludes that if $X$ is an
$S$-prescheme locally of finite type, \emph{$X$ is smooth over~$S$ if and only if it is flat over~$S$ and differentially
smooth over~$S$.} (N.B. the flatness hypothesis is essential,
as one sees by taking for $X$ a closed subprescheme
of~$S$).

Let us note further, by way of recall, that one obtains a
\emph{second algebra structure} on~$\cal{P}_{X/S}^n$
thanks to the projection $\pr_2\colon X\times_S X\to X$,
which is moreover deduced from the former by means of
the \emph{canonical involution} of the sheaf of rings $\cal{P}_{X/S}^n$,
induced by the symmetry automorphism of~$X\times_S X$. One denotes
by $d_{X/S}^n$
or simply $d^n$, the homomorphism of sheaves of rings
\begin{equation}
\label{eq:II.4.6}
d_{X/S}^n\colon \cal{O}_X\to\cal{P}_{X/S}^n
\end{equation}
which corresponds to this second
Algebra structure. Taking the isomorphism~\eqref{eq:II.4.4} into account, this
homomorphism transforms a section $f$ of~$\cal{O}_X$ into a section
$d^n(f)$ of~$\cal{P}_{X/S}^n$ whose inverse image under a section $s$
of~$X$ over~$S$ is identified with the canonical image of~$f$ in
$\Gamma(X,\cal{O}_X/\cal{J}^{n+1})$. This justifies the name
``\emph{system of principal parts of order $n$ of~$f$}''
given to $d^nf$, notably in the case where $S=\Spec(k)$,
% original p. 51
envisaged in formula~\eqref{eq:II.4.5}.

To finish, let us note that the homomorphism \eqref{eq:II.4.6} may be
considered as the \emph{differential operator
of order}~$\leq n$\footnote{For everything concerning the present
paragraph, one may consult EGA~IV~16.8 to 16.12.}
(relative to the prescheme of constants $S$)
\emph{universal}
on~$\cal{O}_X$, by agreeing to call differential operator
of order~$\leq n$
from~$\cal{O}_X$ into a Module $F$, a homomorphism of sheaves $D$
which factors as
\[
D\colon \cal{O}_X\lto{d^n} \cal{P}_{X/S}^n\lto{u} F
\]
where $u$ is a homomorphism \emph{of~$\cal{O}_X$-Modules},
moreover uniquely determined by~$D$. This definition
agrees with the intuitive recursive definition ($D$ is a
differential operator of order $\leq n$ if for every section
$g$ of~$\cal{O}_X$ on an open set $U$ of~$X$, $f\mto D(fg)-D(f)$ is
a differential operator of order $\leq n-1$ on~$U$). It
follows that \emph{if $X$ is differentially smooth over~$S$, the
sheaf of rings of differential operators of all orders
has the simple well-known structure} in differential calculus on
differentiable manifolds, and in particular admits
locally a basis of~$\cal{O}_X$-Module
formed of the \emph{divided powers} in commuting operators
$\delta/\delta x_i$ $(1\leq i\leq n)$. If~$S$ is a sheaf of
$\QQ$-algebras ($\QQ=$ field of rationals) it suffices to take
the ordinary polynomials in the $\delta/\delta x_i$. In this case
moreover, and quite exceptionally, for $X$ to be
differentially smooth over~$S$, it already suffices that
$\it{\Omega}_{X/S}^1$ be locally free.
\end{remarks}
\endgroup

\begin{remark}
\label{II.4.19}
The terminology ``regular immersion'', ``regular
ideal'', etc.\ introduced in this no.\ has met with
a rather lively and general opposition (Chevalley, Serre). One has
proposed ``Cohen--Macaulay ideal'' or ``Macaulay
ideal'' or ``Macaulayan ideal'', which would morally also oblige
one to adopt ``Cohen--Macaulay immersion'' or ``Macaulay
immersion''. This terminology however conflicts with another
already used in future drafts of the
multiplodoque, where a morphism (of finite type) is said to be
``Cohen--Macaulay'' at a point if it is flat at that point, and if the
fiber passing through that point has there a local ring which is a
Cohen--Macaulay ring. Until a satisfactory solution is found,
we shall keep under all reservations the terminology introduced
in this no.\footnote{It is the one adopted in
EGA~$\mathrm{0}_\mathrm{IV}$~15.1.7.}
\end{remark}

\section{The case of a ground field}
\label{II.5}
% original p. 52

\begin{proposition}
\label{II.5.1}
Let $k$ be a field, $X$ a prescheme of finite type over~$k$,
$x$ a point of~$X$ and $n$ the dimension of~$X$ at~$x$,
$f\colon X\to\Spec k[t_1,\dots,t_n]=Y$
a morphism, defined by elements
$f_i\in\Gamma(X,\cal{O}_X)$.
The following conditions are equivalent (and imply that
$X$ is smooth over~$k$ at~$x$, and a fortiori regular at~$x$
by~\ref{II.3.1}):
\begin{enumerate}
\item[(i)] $f$ is \'etale at~$x$.
\item[(ii)] The $df_i$ form a basis of~$\it{\Omega}_{X/k}^1$ at~$x$.
\item[(iii)] The $df_i$ generate $\it{\Omega}_{X/k}^1$ at~$x$.
\end{enumerate}
\end{proposition}

Since (i) implies that $X$ is smooth over~$k$ at~$x$, the implication
(i)$\To$(ii) is a special case of~\ref{II.4.8},
(ii)$\To$(iii) is trivial, it remains to prove
(iii)$\To$(i). Now if (iii) is satisfied, $f$ is net
at~$x$ by virtue of lemma~\ref{II.4.1}, hence (replacing $x$ by
an open neighborhood) quasi-finite, hence dominant for reasons of
dimensions. Since $Y$ is regular, it follows that $f$ is
\'etale by (I~\SourceRef{I.9.5}{9.5}~(ii)) or (I~\SourceRef{I.9.11}{9.11}).

\begin{corollary}
\label{II.5.2}
Under the preliminary conditions of~\ref{II.5.1}, suppose that
$\kres(x)$
is a \emph{finite separable} extension of~$k$, and
that the $f_i$ ($1\leq i\leq n$) define elements
of~$\goth{m}_x$. Then the preceding conditions
are equivalent to
\begin{enumerate}
\item[(iv)] The $f_i$ form a system of generators of
$\goth{m}_x$ (or equivalently: the $f_i$ mod $\goth{m}_x^2$ form a basis
of~$\goth{m}_x/\goth{m}_x^2$
over~$\kres(x)$).
\end{enumerate}
\end{corollary}

Indeed, (iv)$\To$(iii) by virtue of the exact sequence
\begin{equation}
\label{eq:II.5.1}
\goth{m}_x/\goth{m}_x^2\to\Omega_{\cal{O}_x/k}^1\to
\Omega_{\kres(x)/k}^1\to 0
\end{equation}
and taking into account that~$\Omega_{\kres(x)/k}^1=0$ since $\kres(x)$ is
\'etale over~$k$. On the other hand (ii) implies (iv), for since $X$ and
$\Spec(\kres(x))$ are smooth over~$k$ at~$x$, one can in the preceding
exact sequence put a $0$ on the left by virtue of~\ref{II.4.10}~(iv).

\begin{corollary}
\label{II.5.3}
Let $x$ be a point of~$X$ (of finite type over~$k$). If $X$ is
smooth over~$k$ at~$x$, then $\cal{O}_x$ is regular, the converse
being true if
$\kres(x)$
is a finite separable
extension of~$k$.
\end{corollary}

Indeed,
% original p. 53
the converse follows from~\ref{II.5.2}, by taking a
regular system $(f_i)$ of generators of
$\goth{m}_x$. (N.B. instead of~\ref{II.5.2}, one can also invoke
theorem~\ref{II.4.15}). One concludes:

\begin{proposition}
\label{II.5.4}
Let $X$ be a prescheme of finite type over~$k$. If $X$ is smooth
over~$k$, it is regular, the converse being true if $k$
is perfect.
\end{proposition}

For the converse, one notes that by virtue of~\ref{II.5.3}, $X$ is
smooth over~$k$ at every closed point, hence everywhere since the set
of points where it is smooth is open.

\begin{theorem}
\label{II.5.5}
Let $X$ be a prescheme of finite type over~$k$, $x$ a point of
$X$, $n$ the dimension of~$X$ at~$x$, $k'$ a perfect extension of
$k$. The following conditions are equivalent:
\begin{enumerate}
\item[(i)] $X$ is smooth over~$k$ at~$x$.
\item[(ii)] $\it{\Omega}_{X/k}^1$ is free of rank $n$ at~$x$.
\item[(ii~bis)] $\it{\Omega}_{X/k}^1$ is generated by $n$
elements at~$x$.
\item[(iii)] $X$ is differentially smooth over~$k$ at~$x$.
\item[(iv)] There exists an open neighborhood $U$ of~$x$ such that
$U\otimes_k k'$ is regular (\ie the local rings of its
points are regular).
\end{enumerate}
\end{theorem}

One has (i)$\To$(ii) by~\ref{II.4.3}~(ii), (ii)$\To$(ii~bis) trivially and (ii~bis)$\To$(i) thanks
to~\ref{II.5.1}. Since $X$ is flat over~$k$, one has
(i)$\Leftrightarrow$(iii) by virtue of~\ref{II.4.18}. One has
(i)$\To$(iv) since smooth is invariant under extension of the
base and implies regular, and (iv)$\To$(i) for by
Proposition~\ref{II.5.4}, one sees that $U\otimes_k k'$ is simple over~$k'$, hence $U$ is simple over~$k$ by~\ref{II.4.13}.

Taking for $x$ the generic point of~$X$ assumed
irreducible, one finds:

\begin{corollary}
\label{II.5.6}
Let $K$ be a local Artin ring which is a localization of an algebra of
finite type over the field~$k$ (for example, $K$ may be an
extension of finite type of~$k$), let~$n$ be the transcendence
degree over~$K$ of its residue field. Equivalent conditions:
\begin{enumerate}
\item[(i)] $K$ is a finite separable extension of a purely
transcendental extension $k(t_1,\dots,t_n)$ of~$k$.
\item[(ii)]
% original p. 54
$\it{\Omega}^1_{X/k}$
is a $K$-module free of rank~$n$.
\item[(ii~bis)]
$\it{\Omega}^1_{X/k}$
is a $K$-module admitting $n$
generators.
\item[(iii)] The completion $O'$ of~$K\otimes_k K$ for the
topology defined by the powers of the augmentation
ideal $K\otimes_k K\to K$ is a
``regular'' augmented
$K$-algebra,
\ie isomorphic to an
algebra of formal series in~$K$ (N.B. If~$K$ is a field,
this is equivalent to saying that~$O'$ is a
regular local ring).
\item[(iv)] $K$ is a separable extension of~$k$.
\end{enumerate}
\end{corollary}

Indeed, one can always consider $K$ as the local ring of the
generic point of an irreducible scheme of finite type
$X$ over~$k$, and the conditions envisaged are
the conditions of the same name in~\ref{II.5.5}, taking in~(iv)
for~$k'$ an algebraically closed extension of~$k$. Only
the implication $K$ separable over~$k$ $\To$ $X$ smooth over~$k$ at~$x$,
requires a proof. Now one is immediately
reduced thanks to~\ref{II.4.13} to the case where the ground
field is~$k'$, hence algebraically closed, hence where there exists a
point $a$ of~$X$ rational over~$k$. But then $X$ is smooth over~$k$
at~$a$ by~\ref{II.5.4}, a fortiori it is smooth over~$k$ at the
generic point~$x$, qed\footnote{\Cf Errata at the end of the present Exp\ptbl II
(p\ptbl \pageref{II.fin.errata})}.

One will note that in the case where~$K$ is an extension of finite
type of~$k$, the equivalence of
(i), (ii), (ii~bis), (iv)
is well
known, but that we have not made use of any of
these
already known
equivalences.
Of course,
Proposition~\ref{II.5.1} contains as a special case that a sequence
of elements~$x_i$ ($1\leq i\leq n$) is a ``separating
transcendence basis'' of~$K$ over~$k$ if and only if
the~$dx_i$ form a basis of the $K$-module $\Omega^1_{K/k}$, a fact well
known moreover.

\begin{corollary}
\label{II.5.7}
Let $X$ be a prescheme of finite type over a
field~$k$. For $X$ to be smooth over~$k$, it is necessary and sufficient that
$\it{\Omega}^1_{X/k}$
be locally free, and that the local rings of the
generic points of the irreducible components of~$X$
be separable extensions of~$k$ (this last
condition being automatically satisfied if $k$ is
perfect and $X$ reduced).
\end{corollary}

One may assume $X$ connected, let $n$ be the rank of
$\it{\Omega}^1_{X/k}$
assumed locally free. By the hypothesis
and~\ref{II.5.6}, this is also the transcendence degree of the
extensions
% original p. 55
of~$k$ defined by the local rings of the generic
points of~$X$, hence all the
irreducible components of~$X$ are of dimension~$n$. One then concludes
thanks to~\ref{II.5.5}.

One should take care that if $K$ is a finite extension (not
necessarily separable) of~$k$, then $\Omega^1_{K/k}$ is
a free $k$-module, hence setting $X=\Spec(K)$, $\Omega^1_{X/k}$ is a
locally free sheaf, and $X$ is reduced, without $X$ being
necessarily smooth over~$k$. Extending then the scalars to
the algebraic closure of~$k$, one finds an analogous example,
where $k$ is algebraically closed, but $X$ on the other hand
not being reduced.

\begin{corollary}
\label{II.5.8}
Let $X$ be a prescheme of finite type over the field~$k$, $x$ a
point of~$X$, $n$ the dimension of~$X$ at~$x$, $p$ the dimension
of~$\cal{O}_x$,
\ie the codimension in $X$ of the closure $Y$ of
$x$ in~$X$; hence $n-p$ the transcendence degree
of~$\kres(x)$
over~$k$. Let $f_i$ ($1\leq i\leq n$) be elements
of~$\cal{O}_x$, such that $f_i\in \goth{m}_x$ for $1\leq i\leq p$. The
following conditions are equivalent

\begin{enumerate}
\item[(i)] the germ of morphism at~$x$
\[
X\to \Spec\bigl(k[t_1,\dots,t_n]\bigr)
\]
defined by the $f_i$ is \'etale at~$x$.
\item[(ii)] The $f_i$ ($1\leq i\leq p$)
generate~$\goth{m}_x$,
\ie form a regular system of parameters
of~$\cal{O}_x$, and the classes in
$\kres(x)$ of the~$f_j$ ($p+1\leq j\leq n$) form a separating
transcendence basis, (\ie the~$d\overline{f}_j$ ($p+1\leq j\leq n$) form a basis
of~$\Omega^1_{\kres(x)/k}$, or equivalently generate~$\Omega^1_{\kres(x)/k}$).
\end{enumerate}
\end{corollary}

Suppose~(i) is satisfied. It follows that the $df_i(x)$
form a basis of~$\it{\Omega}^1_{X/k}(x)$ \eqref{II.4.8} hence their
images $d\overline{f}_i(x)$ in $\Omega^1_{\kres(x)/k}$ generate
this vector space over~$k$. Since the $\overline{f}_i$ for $1\leq
i\leq p$ are zero, it follows that it suffices to take the
$d\overline{f}_i(x)$ with $p+1\leq i\leq n$. Since the
transcendence degree of~$\kres(x)$ over~$k$ is~$n-p$, it then follows
from corollary~\ref{II.5.6} criterion~(iii) (applied
to~$K=\kres(x)$) that $Y$ is smooth over~$k$ at its
generic point~$x$, and that the~$d\overline{f}_i(x)$ ($p+1\leq i\leq
n$) form a \emph{basis} of~$\Omega^1_{\kres(x)/k}$
over~$\kres(x)$. Consequently, condition~(ii) of~\ref{II.4.9} is
satisfied, hence also condition~(iii) and in particular
the~$f_i$ ($1\leq i\leq p$) form a system of
generators of~$\goth{m}_x$. Since
% original p. 56
$\cal{O}_x$ is of dimension~$p$, they therefore form a
regular system of parameters at~$x$. This proves~(ii).

Suppose~(ii) is satisfied. By virtue of the exact
sequence~\eqref{II.5.1}, it follows that the~$df_i(x)$
generate~$\it{\Omega}^1_{X/k}$,
whence~(i) thanks to prop\ptbl\ref{II.5.1}.

\begin{corollary}
\label{II.5.9}
Let~$X$ be a prescheme of finite type over the
field~$k$, $x$ a point of~$X$, $n$ the dimension of~$X$ at~$x$, $p$ the
dimension
of~$\cal{O}_x$,
\ie the codimension of the closure~$Y$
of~$x$ in~$X$, hence~$n-p$ the transcendence degree
of~$\kres(x)$
over~$k$. Equivalent conditions:
\begin{enumerate}
\item[(i)] $\cal{O}_x$ is regular
and~$\kres(x)$
 is a
separable extension of~$k$.
\item[(ii)] $X$ is smooth over~$k$ at~$x$, and the canonical homomorphism
\[
\goth{m}_x/\goth{m}_x^2\to
\Omega^1_{\cal{O}_x/k}\otimes_{\cal{O}_x}
\kres(x)=\it{\Omega}^1_{X/k}(x)
\]
is injective.
\item[(iii)] There are~$f_i\in \cal{O}_x$ ($1\leq i\leq n$)
with~$f_i\in \goth{m}_x$ for~$1\leq i\leq p$, such that the germ of
morphism at~$x$ from~$X$ to~$\Spec\bigl(k[t_1,\dots,t_n]\bigr)$
defined by the~$f_i$ is \'etale at~$x$ (\ie by~\ref{II.5.1}
such that the~$df_i(x)$ generate~$\it{\Omega}^1_{X/k}(x)$).
\item[(iv)] There are~$f_i\in \cal{O}_x$ ($1\leq i\leq n$) such that
the~$f_i$ ($1\leq i\leq p$) generate~$\goth{m}_x$ and that
the~$df_j(x)$ ($p+1\leq j\leq n$)
generate~$\Omega^1_{\kres(x)/k}$ over~$\kres(x)$.
\end{enumerate}
\end{corollary}

The equivalence of~(iii) and~(iv) follows from
corollary~\ref{II.5.8}, these conditions are also equivalent
by~\ref{II.4.9} to the fact that~$X$ is smooth over~$k$ at~$x$ and
that condition~(ii)
of~\ref{II.4.10}
is satisfied. Hence they are equivalent to the fact that~$X$ is
smooth over~$k$ at~$x$ and that condition~(iv) of~\ref{II.4.10} is
satisfied, hence to~\ref{II.5.9}~(ii). Or to the fact that~$X$ is
smooth over~$k$ at~$x$ and that condition~(i) of~\ref{II.4.10} is
satisfied, which here simply means
that~$\kres(x)$
is
separable over~$k$. This implies~\ref{II.5.9}~(i), it remains to
prove that~\ref{II.5.9}~(i) implies it, \ie to prove the

\begin{corollary}
\label{II.5.10}
Let~$x$ be a point
of a
prescheme of finite type over the field~$k$, such
that~$\kres(x)$
is separable over~$k$. For~$X$ to be smooth over~$k$ at~$x$,
it is necessary and sufficient that it be regular at~$x$ (\ie that the local
ring~$\cal{O}_x$ be regular).
\end{corollary}

Indeed, if this is so, one can evidently find
$f_i\in \cal{O}_x$ ($1\leq i\leq n$) satisfying
condition~\ref{II.5.9}~(iv).

\subsection*{Errata}
\label{II.fin.errata}
In
% original p. 57
the present no., proof of~\ref{II.5.6}, one
has used the fact that a nonempty reduced scheme of finite type
over an algebraically closed field admits at least one
regular point (hence smooth), a fact which is usually proved by
the differential method (via Zariski's theorem that
the set of regular points of~$X$ is open). If one wants
to avoid a vicious circle, one must prove that if~$K/k$ is
a separable extension of finite type, and if the~$f_i\in K$ are
such that the~$d_{K/k}f_i$ form a basis of~$\Omega^1_{K/k}$, ($1\leq
i\leq n$), then~$n$ is the transcendence degree of~$K$
over~$k$,
\ie the~$f_i$ are algebraically independent. The
proof of this fact with the aid of the
Mac\kern2pt Lane
criterion is well known, \cf Bourbaki, Alg\`ebre, Chap\ptbl V par\ptbl 9 th\ptbl 2: one
takes a polynomial~$g\in k[t_1,\dots,t_n]$ of minimal
degree such that~$g(f_1,\dots,f_n)=0$, one has therefore
\[
\sum\frac{dg}{d t_i}(f_1,\dots,f_n) df_i=0
\]
whence (since the~$df_i$ form a basis of~$\Omega^1_{K/k}$) the
fact that the~$dg/d t_i$ vanish at~$(f_1,\dots,f_n)$, hence are zero
by the minimal character of~$g$, hence if~$k$ is of
characteristic~$0$ one has~$g=0$, and if~$k$ is of
characteristic~$p\neq 0$ one has~$g=h(t_1^p,\dots,t_n^p)$.
Using the
Mac\kern2pt Lane
criterion,
one sees that the
polynomial~$h\in k[t_1,\dots,t_n]$ also vanishes at~$(f_1,\dots,f_n)$,
whence again~$g=0$ by the minimal character of~$g$.

\end{document}
